% Options for packages loaded elsewhere
\PassOptionsToPackage{unicode}{hyperref}
\PassOptionsToPackage{hyphens}{url}
%
\documentclass[
  10pt,
  a4paper,
]{article}
\usepackage{amsmath,amssymb}
\usepackage{iftex}
\ifPDFTeX
  \usepackage[T1]{fontenc}
  \usepackage[utf8]{inputenc}
  \usepackage{textcomp} % provide euro and other symbols
\else % if luatex or xetex
  \usepackage{unicode-math} % this also loads fontspec
  \defaultfontfeatures{Scale=MatchLowercase}
  \defaultfontfeatures[\rmfamily]{Ligatures=TeX,Scale=1}
\fi
\usepackage{lmodern}
\ifPDFTeX\else
  % xetex/luatex font selection
\fi
% Use upquote if available, for straight quotes in verbatim environments
\IfFileExists{upquote.sty}{\usepackage{upquote}}{}
\IfFileExists{microtype.sty}{% use microtype if available
  \usepackage[]{microtype}
  \UseMicrotypeSet[protrusion]{basicmath} % disable protrusion for tt fonts
}{}
\makeatletter
\@ifundefined{KOMAClassName}{% if non-KOMA class
  \IfFileExists{parskip.sty}{%
    \usepackage{parskip}
  }{% else
    \setlength{\parindent}{0pt}
    \setlength{\parskip}{6pt plus 2pt minus 1pt}}
}{% if KOMA class
  \KOMAoptions{parskip=half}}
\makeatother
\usepackage{xcolor}
\usepackage[margin=22mm]{geometry}
\setlength{\emergencystretch}{3em} % prevent overfull lines
\providecommand{\tightlist}{%
  \setlength{\itemsep}{0pt}\setlength{\parskip}{0pt}}
\setcounter{secnumdepth}{-\maxdimen} % remove section numbering
\ifLuaTeX
\usepackage[bidi=basic]{babel}
\else
\usepackage[bidi=default]{babel}
\fi
\babelprovide[main,import]{italian}
% get rid of language-specific shorthands (see #6817):
\let\LanguageShortHands\languageshorthands
\def\languageshorthands#1{}
\ifLuaTeX
  \usepackage{selnolig}  % disable illegal ligatures
\fi
\usepackage{bookmark}
\IfFileExists{xurl.sty}{\usepackage{xurl}}{} % add URL line breaks if available
\urlstyle{same}
\hypersetup{
  pdflang={it},
  hidelinks,
  pdfcreator={LaTeX via pandoc}}

\author{}
\date{}

\begin{document}

{
\setcounter{tocdepth}{3}
\tableofcontents
}
\section{ThreatForge MR comparison - current vs
proposed}\label{threatforge-mr-comparison---current-vs-proposed}

\begin{quote}
Non-canonical case study. Baseline for current documents is the pinned
ThreatForge product-semantic commit below.

\texttt{cae0f7b6b37f430ac4e857aabf6ef9f87c89dbb1}
\end{quote}

The purpose is to test the general MR metamodel against a real project
and to collect refactoring/authoring evidence. The proposed rewrites are
\textbf{not canonical ThreatForge changes}.

\subsection{MR-0001}\label{mr-0001}

\subsubsection{Current document}\label{current-document}

\paragraph{Original: MR-0001 - Gestione documentale
governata}\label{original-mr-0001---gestione-documentale-governata}

\subparagraph{Intent}\label{intent}

Define a governed documentation system that gives ThreatForge one
concise, analyzable and navigable source for project intent, decisions,
requirements, assets, relations and verification evidence.

\subparagraph{Context}\label{context}

ThreatForge treats governed documentation as the primary project model
for people, developers, LLM consumers and deterministic checks.
Canonical registries and Markdown bodies provide the maintained sources,
while readable books, HTML views, diagrams, indexes, appendices and
reports are derived projections.

The model organizes project knowledge by Macro-requirement and uses
explicit Decisions and Requirements to refine each project theme.
Controlled vocabularies, taxonomies, stable identifiers and explicit
relations reduce ambiguity and support future security and threat
analysis.

\subparagraph{Macro obligation}\label{macro-obligation}

\begin{itemize}
\tightlist
\item
  ThreatForge must organize governed project knowledge around stable
  Macro-requirement identities.
\item
  Each Macro-requirement must provide the project-level context for its
  owned Decisions and Requirements.
\item
  Governed documentation must separate canonical sources from derived
  readable projections.
\item
  Canonical registries must support both navigation and deterministic
  validation.
\item
  Governed Markdown bodies must remain concise, analyzable and linked to
  their canonical registry records.
\item
  Governed documents must use stable identifiers and canonical
  terminology.
\item
  Controlled documentation fields must reference governed value sets,
  taxonomies or vocabularies.
\item
  Governed documentation must distinguish authored prose, required
  fields, controlled fields and derived values.
\item
  Decision records must govern specific choices within one
  Macro-requirement.
\item
  Functional Requirements must describe independently testable
  capabilities, behaviors or outcomes.
\item
  Specialized Requirements must attach governed constraints, controls or
  qualities to Functional Requirements.
\item
  Governed assets must use stable identities and canonical descriptions
  when they enter the project model.
\item
  Governed relations must connect Macro-requirements, Decisions,
  Requirements, assets, implementations and verification evidence
  explicitly.
\item
  ThreatForge must organize tutorials, how-to guides, reference material
  and explanations according to Diátaxis.
\item
  Explanation and how-to documents must not become independent canonical
  sources of governed obligations.
\item
  Derived documentation must obtain indexes, linked records and metadata
  from canonical sources instead of manual duplication.
\item
  Governed documentation must preserve sufficient lifecycle and
  historical traceability for superseded, deprecated or removed records.
\item
  Deterministic corpus analysis must identify terminology drift,
  ambiguous labels and uncontrolled references.
\item
  Governed documentation must support future graph exploration and
  threat-analysis workflows without duplicating canonical facts.
\end{itemize}

\subparagraph{Scope}\label{scope}

\begin{itemize}
\tightlist
\item
  Includes: Macro-requirement registry records and Markdown bodies
\item
  Includes: Decision and Requirement registries and bodies
\item
  Includes: {[}BAE-0001{]} Governed project documentation
\item
  Includes: governed assets, controlled vocabularies and taxonomies
\item
  Includes: stable identifiers, controlled fields and explicit relations
\item
  Includes: implementation and verification traceability
\item
  Includes: Diátaxis documentation connected to governed sources
\item
  Includes: deterministic corpus-quality reports and derived
  documentation projections
\item
  Excludes: final application backend and user-interface implementation
\item
  Excludes: runtime threat-analysis execution
\item
  Excludes: autonomous AI integrations and model runtime behavior
\item
  Excludes: ungoverned narrative content without a governed
  project-model relation
\end{itemize}

\subparagraph{Non-goals}\label{non-goals}

\begin{itemize}
\tightlist
\item
  Mandatory adoption of RDF, SKOS, SHACL or OWL in the initial model
\item
  Manual maintenance of indexes or metadata already derivable from
  canonical sources
\item
  Use of explanation or how-to documents as autonomous normative sources
\end{itemize}

\subsubsection{Proposed rewrite for metamodel
test}\label{proposed-rewrite-for-metamodel-test}

\paragraph{Proposed: MR-0001 - Documentazione di progetto governata e
tracciabile}\label{proposed-mr-0001---documentazione-di-progetto-governata-e-tracciabile}

\subparagraph{Intent}\label{intent-1}

Consentire a persone e strumenti di comprendere, mantenere e verificare
le conoscenze che guidano un progetto attraverso documentazione
versionata, coerente e tracciabile.

\subparagraph{Context}\label{context-1}

ThreatForge usa la documentazione governata come base condivisa per
descrivere intenzioni, scelte, requisiti e relative evidenze. La
documentazione deve rimanere leggibile dalle persone e sufficientemente
strutturata da poter essere verificata e navigata dagli strumenti.

\subparagraph{Stakeholders}\label{stakeholders}

\begin{itemize}
\tightlist
\item
  responsabili di progetto
\item
  sviluppatori e manutentori
\item
  reviewer
\item
  analisti
\item
  utenti che governano progetti con ThreatForge
\end{itemize}

\subparagraph{Objective}\label{objective}

\begin{itemize}
\tightlist
\item
  Rendere la documentazione una fonte affidabile per comprendere e
  governare l'evoluzione del progetto.
\item
  Mantenere espliciti i collegamenti tra intenzioni, decisioni,
  requisiti ed evidenze senza duplicare le stesse regole in piu livelli.
\end{itemize}

\subparagraph{Scope (candidate semantics under
test)}\label{scope-candidate-semantics-under-test}

\begin{itemize}
\tightlist
\item
  conoscenza governata del progetto e sua tracciabilita
\item
  leggibilita umana e verificabilita automatizzabile della
  documentazione
\end{itemize}

\subsubsection{What moved down}\label{what-moved-down}

Registri, Markdown body, Diataxis, materializzazioni, tassonomie,
controlli terminologici e meccanismi di proiezione diventano
Decision/Requirement/Governance detail.

\begin{center}\rule{0.5\linewidth}{0.5pt}\end{center}

\subsection{MR-0002}\label{mr-0002}

\subsubsection{Current document}\label{current-document-1}

\paragraph{Original: MR-0002 - Authoring e implementazione
governata}\label{original-mr-0002---authoring-e-implementazione-governata}

\subparagraph{Intent}\label{intent-2}

Define a governed path from an existing Requirement to planned, created
and verified implementation artifacts through deterministic core
capabilities and thin development-environment adapters.

\subparagraph{Context}\label{context-2}

ThreatForge supports people and LLM consumers while they author governed
Requirements and connect those Requirements to implementation and
verification evidence. The core workflow remains independent from a
specific editor, while Visual Studio Code is the first integrated
authoring surface.

The repository runner, authoring tools, implementation planners,
scaffolders and promotion controls form one governed workflow. Editor
integrations collect input, invoke those capabilities and present their
results without owning domain rules or bypassing repository gates.

\subparagraph{Macro obligation}\label{macro-obligation-1}

\begin{itemize}
\tightlist
\item
  Every implementation artifact must derive from at least one existing
  governed Requirement.
\item
  ThreatForge must keep domain rules and traceability rules inside
  governed core modules and tools.
\item
  Development-environment adapters must remain thin consumers of
  governed capabilities.
\item
  Requirement authoring must use canonical registries, controlled value
  sets and deterministic body profiles.
\item
  Requirement creation must provide a complete preview before persistent
  repository changes.
\item
  Persistent authoring operations must require explicit confirmation.
\item
  Governed creation operations must preserve atomicity across registry
  and body updates.
\item
  Governed creation operations must restore pre-operation state after
  verification failure.
\item
  Implementation planning must remain read-only before artifact
  creation.
\item
  Every implementation plan must identify the governing Requirement.
\item
  Implementation scaffolding must create traceable artifacts with
  controlled lifecycle state.
\item
  Implementation promotion must require completed source content and
  successful verification evidence.
\item
  Implementation artifacts must declare bidirectional traceability to
  their linked Requirements.
\item
  Repository commit and push operations must execute registered
  materializers and read-only gates before Git staging.
\item
  Repository projection materialization must remain deterministic,
  bounded and idempotent.
\item
  Direct ungoverned Git operations must not replace the governed
  repository runner.
\item
  Visual Studio Code tasks must delegate to governed commands without
  duplicating canonical choices or validation logic.
\item
  Future editor adapters must consume the same governed capabilities
  without changing the canonical core rules.
\item
  The ThreatForge application must orchestrate the complete workflow
  without coupling the model to one editor or filesystem adapter.
\end{itemize}

\subparagraph{Scope}\label{scope-1}

\begin{itemize}
\tightlist
\item
  Includes: governed Requirement authoring and deterministic previews
\item
  Includes: implementation planning, scaffolding and promotion
\item
  Includes: traceability between Requirements, source artifacts and
  verification evidence
\item
  Includes: registered repository projections and governed commit-push
  operations
\item
  Includes: thin Visual Studio Code integration
\item
  Includes: future editor and application adapters over the same core
  capabilities
\item
  Excludes: autonomous code generation without an existing governed
  Requirement
\item
  Excludes: domain rules duplicated inside editor extensions
\item
  Excludes: ungoverned automatic commit or push behavior
\item
  Excludes: simultaneous initial support for every development
  environment
\end{itemize}

\subparagraph{Non-goals}\label{non-goals-1}

\begin{itemize}
\tightlist
\item
  Complete ThreatForge application user interface in the initial
  authoring workflow
\item
  Editor-specific ownership of canonical validation or controlled values
\item
  Creation of implementation artifacts without governed planning and
  traceability
\end{itemize}

\subsubsection{Proposed rewrite for metamodel
test}\label{proposed-rewrite-for-metamodel-test-1}

\paragraph{Proposed: MR-0002 - Sviluppo tracciabile dai
requisiti}\label{proposed-mr-0002---sviluppo-tracciabile-dai-requisiti}

\subparagraph{Intent}\label{intent-3}

Consentire al lavoro di implementazione di partire da requisiti
governati e di mantenere chiaro, durante l'evoluzione del progetto, che
cosa viene realizzato e con quali evidenze viene verificato.

\subparagraph{Context}\label{context-3}

ThreatForge deve aiutare chi sviluppa un progetto a passare dalla
documentazione governata al lavoro implementativo senza perdere il
legame con le ragioni e gli obblighi che hanno originato quel lavoro.

\subparagraph{Stakeholders}\label{stakeholders-1}

\begin{itemize}
\tightlist
\item
  sviluppatori
\item
  reviewer
\item
  responsabili di progetto
\item
  manutentori
\item
  utenti che usano ThreatForge per guidare l'authoring e
  l'implementazione
\end{itemize}

\subparagraph{Objective}\label{objective-1}

\begin{itemize}
\tightlist
\item
  Mantenere tracciabile il percorso dal requisito al lavoro
  implementativo e alla verifica.
\item
  Offrire assistenza coerente all'authoring e all'esecuzione del lavoro
  senza spostare le regole di dominio dentro uno specifico strumento di
  sviluppo.
\end{itemize}

\subparagraph{Scope (candidate semantics under
test)}\label{scope-candidate-semantics-under-test-1}

\begin{itemize}
\tightlist
\item
  passaggio governato da requisiti a implementazione e verifica
\item
  tracciabilita dello stato del lavoro rispetto ai requisiti
\end{itemize}

\subsubsection{What moved down}\label{what-moved-down-1}

VS Code, repository runner, preview, atomicita, rollback, scaffolding,
promotion, Git staging e materializer sono scelte/obblighi di livello
inferiore.

\begin{center}\rule{0.5\linewidth}{0.5pt}\end{center}

\subsection{MR-0003}\label{mr-0003}

\subsubsection{Current document}\label{current-document-2}

\paragraph{Original: MR-0003 - Modello di Analisi Base derivato dalla
documentazione}\label{original-mr-0003---modello-di-analisi-base-derivato-dalla-documentazione}

\subparagraph{Intent}\label{intent-4}

Define a methodology-neutral and documentation-derived system model that
transforms governed project knowledge into a canonical inventory of Base
Analysis Elements for threat analysis before implementation and
throughout system evolution.

\subparagraph{Context}\label{context-4}

MR-0001 establishes governed documentation as the primary project model,
while MR-0002 governs document authoring and the path from Requirements
to implementation artifacts.

ThreatForge currently lacks an explicit analytical representation
derived from governed documentation without replacing its canonical
facts. Base Analysis Elements provide stable identities, provenance and
relations for reconstructing the system being designed, identifying
incomplete knowledge and determining whether prior analyses are stale.

This boundary supports moving threat analysis into the documentation and
design phase instead of waiting for implementation artifacts to exist.

\subparagraph{Macro obligation}\label{macro-obligation-2}

\begin{itemize}
\tightlist
\item
  ThreatForge must maintain one canonical inventory of Base Analysis
  Elements for each analyzed project.
\item
  Every Base Analysis Element must have a stable governed identity,
  canonical type, canonical title, explicit meaning, provenance and
  lifecycle state.
\item
  Every Base Analysis Element must be justified by existing governed
  project knowledge or by an explicit reviewed analytical addition.
\item
  The governed origin of a Base Analysis Element must exist before any
  document that references that element.
\item
  A document must not reference a Base Analysis Element whose governed
  origin belongs to one of that document's descendants.
\item
  Base Analysis Element provenance must identify the governed documents
  and evidence from which the element is derived.
\item
  The Base Analysis model must preserve explicit relations among
  elements and between elements, Decisions, Requirements and other
  governed project records.
\item
  The Base Analysis inventory must remain methodology-neutral.
\item
  The Base Analysis inventory must not contain STRIDE, STRIDE-AI or
  other overlay-specific classifications as base facts.
\item
  Threat-analysis overlays must consume the canonical Base Analysis
  inventory without silently adding, removing or redefining its
  elements.
\item
  Missing or changed system knowledge discovered by an overlay must
  produce a governed proposal to update the Base Analysis inventory.
\item
  Changes to governed source documentation must support deterministic
  identification of Base Analysis Elements and analyses that are
  potentially stale.
\item
  The Base Analysis model must support analysis before implementation
  and re-analysis after architectural, requirement or feature changes.
\item
  MR-0003 must own Base Analysis Element semantics, provenance,
  relations and lifecycle without owning the generic Markdown reference
  syntax governed by MR-0001.
\item
  MR-0003 must expose canonical Base Analysis projections that MR-0002
  authoring tools and editor adapters can consume without duplicating
  domain rules.
\end{itemize}

\subparagraph{Scope}\label{scope-2}

\begin{itemize}
\tightlist
\item
  Includes: Methodology-neutral Base Analysis model
\item
  Includes: Canonical Base Analysis Element inventory
\item
  Includes: Stable Base Analysis Element identities and canonical types
\item
  Includes: Documentary provenance and governed origin
\item
  Includes: Relations among Base Analysis Elements and governed project
  records
\item
  Includes: Base Analysis Element lifecycle and staleness
\item
  Includes: Derivation from Macro-requirements, Decisions, Requirements
  and other governed project documentation
\item
  Includes: Reviewed analytical additions when governed documentation is
  incomplete
\item
  Includes: Canonical system topology and data-flow projections for
  later threat-analysis overlays
\item
  Includes: Traceability from source documentation to Base Analysis
  Elements and dependent analyses
\item
  Excludes: Generic Markdown syntax for governed entity references
\item
  Excludes: Editor-owned completion, hover, diagnostics or quick-fix
  rules
\item
  Excludes: STRIDE and STRIDE-AI classifications
\item
  Excludes: Methodology-specific threats, findings, mitigations or
  security requirements
\item
  Excludes: Runtime Base Analysis implementation and user-interface
  behavior
\item
  Excludes: Automatic canonical acceptance of LLM-inferred elements
\item
  Excludes: Replacement of Macro-requirement, Decision or Requirement
  registries
\end{itemize}

\subparagraph{Non-goals}\label{non-goals-2}

\begin{itemize}
\tightlist
\item
  Freeze the complete Base Analysis Element taxonomy in the
  Macro-requirement
\item
  Select the physical registry or database schema
\item
  Implement extraction, validation, graph or editor tooling
\item
  Treat the first textual occurrence of an element as its governed
  origin
\item
  Make an LLM inference canonical without explicit governed review
\item
  Duplicate project facts that already have an authoritative governed
  source
\end{itemize}

\subsubsection{Proposed rewrite for metamodel
test}\label{proposed-rewrite-for-metamodel-test-2}

\paragraph{Proposed: MR-0003 - Modello del progetto per
l'analisi}\label{proposed-mr-0003---modello-del-progetto-per-lanalisi}

\subparagraph{Intent}\label{intent-5}

Rendere la conoscenza governata del progetto utilizzabile per analisi
diverse anche prima che esista una implementazione completa, mantenendo
chiara la provenienza delle informazioni e gli effetti delle modifiche.

\subparagraph{Context}\label{context-5}

Le analisi hanno bisogno di una rappresentazione coerente di cio che il
progetto descrive. Questa rappresentazione deve derivare dalla
documentazione senza sostituirne i fatti canonici e deve poter evolvere
quando cambia il progetto.

\subparagraph{Stakeholders}\label{stakeholders-2}

\begin{itemize}
\tightlist
\item
  analisti
\item
  architetti e progettisti
\item
  sviluppatori
\item
  reviewer
\item
  responsabili di progetto
\end{itemize}

\subparagraph{Objective}\label{objective-2}

\begin{itemize}
\tightlist
\item
  Derivare dalla documentazione una base comune e tracciabile per
  l'analisi.
\item
  Permettere analisi anticipate e ripetute durante l'evoluzione del
  progetto.
\item
  Mantenere distinguibili i fatti del progetto dalle interpretazioni
  introdotte dalle singole analisi.
\end{itemize}

\subparagraph{Scope (candidate semantics under
test)}\label{scope-candidate-semantics-under-test-2}

\begin{itemize}
\tightlist
\item
  conoscenza del progetto necessaria a costruire una base comune per
  analisi
\item
  provenienza e aggiornamento della conoscenza analitica comune
\end{itemize}

\subsubsection{What moved down in the first
reformulation}\label{what-moved-down-in-the-first-reformulation}

BAE fields, precise provenance rules, topology/data-flow projection,
STRIDE exclusions and editor integration move below MR.

\subsubsection{Second reformulation after ADR-family
review}\label{second-reformulation-after-adr-family-review}

\paragraph{Current candidate: MR-0003 - Governed representation of the
system}\label{current-candidate-mr-0003---governed-representation-of-the-system}

\subparagraph{Intent}\label{intent-6}

Make the system described by the project understandable through
explicit, traceable and maintainable project knowledge, so that people
and downstream activities can rely on a shared representation of what
the system is and how that knowledge is justified.

\subparagraph{Context}\label{context-6}

Project documentation may describe relevant parts of a system across
different documents and at different levels of detail. Some knowledge is
stated directly, while other knowledge may need to be reconstructed,
clarified or added through review.

Without a governed representation, different readers or downstream
activities may interpret the same project differently, create competing
identities for the same concepts, lose the origin of important
knowledge, or continue relying on information that is no longer current.

\subparagraph{Stakeholders}\label{stakeholders-3}

\begin{itemize}
\tightlist
\item
  Project owners and maintainers
\item
  Architects and developers
\item
  Reviewers
\item
  Analysts who consume project knowledge
\item
  Other stakeholders who need a dependable understanding of the
  documented system
\end{itemize}

\subparagraph{Objective}\label{objective-3}

Provide one coherent and traceable representation of the system
described by the project, preserving the distinction between documented
knowledge, reviewed additions and later interpretations.

\subparagraph{Scope}\label{scope-3}

Includes the project knowledge required to identify and understand the
relevant parts of the system and their relationships.

Includes the origin, supporting evidence and evolution of that knowledge
when these are necessary to determine whether it can still be relied
upon.

Includes a governed way to review knowledge that is inferred,
reconstructed or found to be incomplete.

Excludes interpretations whose meaning depends on a particular analysis
method.

Excludes findings, threats, risk judgments and other conclusions
produced by applying an analysis method.

\subparagraph{Assumptions}\label{assumptions}

Project documentation may be incomplete and may evolve over time.

Some relevant system knowledge may require explicit review before it can
be treated as governed project knowledge.

\subparagraph{Constraints}\label{constraints}

The governed representation must remain independent from the terminology
and classifications of any particular analysis method.

Knowledge introduced through inference or analytical discovery must not
become governed project knowledge without explicit review.

\subsubsection{Why the first reformulation was
reopened}\label{why-the-first-reformulation-was-reopened}

The first reformulation described MR-0003 mainly through its future use
by analysis (``a common base for analysis''). Read beside the first
MR-0005 rewrite, this made the two MRs appear almost equivalent.

Before merging them, the historical Decision family was inspected. The
MR-0003 Decisions consistently govern a different concern:
representation of system knowledge, identity and provenance, continuity
of source authority, documentary extraction/review, and the boundary
between system facts and methodological interpretation.

The second reformulation therefore expresses the value of MR-0003
independently from its analytical consumer: \textbf{what do we know
about the system, and why can that knowledge be relied upon?}

\begin{center}\rule{0.5\linewidth}{0.5pt}\end{center}

\subsection{MR-0004}\label{mr-0004}

\subsubsection{Current document}\label{current-document-3}

\paragraph{Original: MR-0004 - Ciclo di vita governato dei progetti
target}\label{original-mr-0004---ciclo-di-vita-governato-dei-progetti-target}

\subparagraph{Intent}\label{intent-7}

Define how ThreatForge creates, accesses, validates, isolates and
analyzes governed Target Projects at user-selected filesystem locations.

\subparagraph{Context}\label{context-7}

ThreatForge needs one small and reproducible path from project creation
to governed documentation and documentation-derived Base Analysis. An
internal demonstration project and a newly created external project use
the same Target Project model even though their destination locations
differ. The ThreatForge engine remains separate from every target, and
the initial demonstrable target can represent a system through governed
documentation before executable application code exists.

\subparagraph{Macro obligation}\label{macro-obligation-3}

\begin{itemize}
\tightlist
\item
  ThreatForge must use one Target Project model for an internal
  demonstration project and a newly created external project.
\item
  Target Project creation must receive one explicit destination root
  selected for the creation request.
\item
  Subsequent Target Project authoring, verification and analysis must
  receive one explicit target root.
\item
  Internal and external Target Projects must originate from the same
  governed template and application behavior.
\item
  Every Target Project must own its project-local documentation,
  registries, reports and materialized projections.
\item
  A Target Project must remain structurally valid without executable
  source code, a running backend, a frontend or a database instance.
\item
  Every Target Project must support the governed document models
  required by the active engine-owned target template and project-local
  Base Analysis records.
\item
  Target Project documentation must provide project-local sources for
  actors, logical components, data resources, trust boundaries and
  information flows.
\item
  Target Project records must preserve stable project-local identifiers
  and source paths.
\item
  Base Analysis Element provenance must resolve to governed sources
  owned by the analyzed Target Project or to explicit reviewed
  analytical additions.
\item
  Target Project creation and use must not modify or contaminate the
  canonical ThreatForge project model.
\item
  ThreatForge repository verification must remain development
  governance.
\item
  ThreatForge repository verification must not be represented as a
  Target Project kind.
\item
  MR-0001 must remain authoritative for governed documentation
  structure, document semantics and Diátaxis organization.
\item
  MR-0002 must remain authoritative for reusable application interfaces,
  target-access boundaries and delivery-adapter separation.
\item
  MR-0003 must remain authoritative for Base Analysis Element semantics,
  provenance, relations and lifecycle.
\item
  Target-specific product requirements must remain inside the governed
  project model of the owning Target Project.
\end{itemize}

\subparagraph{Scope}\label{scope-4}

\begin{itemize}
\tightlist
\item
  Includes: One Target Project model for internal and external
  destinations
\item
  Includes: Explicit destination-root project creation
\item
  Includes: Explicit target-root authoring, verification and analysis
\item
  Includes: Shared governed Target Project template
\item
  Includes: Project-local governed documentation and registries
\item
  Includes: Document-only Target Projects without executable application
  code
\item
  Includes: Project-local sources for actors, components, data
  resources, boundaries and flows
\item
  Includes: Base Analysis readiness and documentary provenance
\item
  Includes: Isolation of target records, reports and materializations
\item
  Includes: Minimal Target Project lifecycle orchestration
\item
  Excludes: Import or migration of arbitrary existing repositories
\item
  Excludes: Compatibility-version negotiation or migration mechanisms
\item
  Excludes: Multiple concurrent target sessions
\item
  Excludes: Final web-interface behavior
\item
  Excludes: Product-domain requirements of a specific Target Project
\item
  Excludes: STRIDE, STRIDE-AI or other methodology overlays
\item
  Excludes: Executable application code as a prerequisite for project
  analysis
\end{itemize}

\subparagraph{Non-goals}\label{non-goals-3}

\begin{itemize}
\tightlist
\item
  Model ThreatForge repository verification as a Target Project kind
\item
  Define compatibility versions, gate profiles or migration mechanisms
\item
  Import arbitrary unstructured repositories
\item
  Implement multiple concurrent Target Project sessions
\item
  Define the final web user interface
\item
  Define the product-domain requirements of the first demonstration
  Target Project
\item
  Require executable application code before documentation-derived Base
  Analysis
\end{itemize}

\subsubsection{Proposed rewrite for metamodel
test}\label{proposed-rewrite-for-metamodel-test-3}

\paragraph{Proposed: MR-0004 - Gestione isolata dei progetti
governati}\label{proposed-mr-0004---gestione-isolata-dei-progetti-governati}

\subparagraph{Intent}\label{intent-8}

Consentire a ThreatForge di creare, aprire e governare progetti distinti
mantenendo separati i loro documenti, risultati e ciclo di vita dal
motore ThreatForge e dagli altri progetti.

\subparagraph{Context}\label{context-8}

ThreatForge e uno strumento riutilizzabile su piu progetti. Ogni
progetto governato deve poter possedere la propria conoscenza e i propri
risultati senza dipendere dal repository del motore o contaminare altri
progetti.

\subparagraph{Stakeholders}\label{stakeholders-4}

\begin{itemize}
\tightlist
\item
  utenti di ThreatForge
\item
  team dei progetti governati
\item
  manutentori di ThreatForge
\item
  reviewer e analisti che lavorano su un progetto specifico
\end{itemize}

\subparagraph{Objective}\label{objective-4}

\begin{itemize}
\tightlist
\item
  Fornire a ogni progetto governato uno spazio autonomo e
  identificabile.
\item
  Mantenere separata l'evoluzione del progetto governato dall'evoluzione
  del motore ThreatForge.
\end{itemize}

\subparagraph{Scope (candidate semantics under
test)}\label{scope-candidate-semantics-under-test-3}

\begin{itemize}
\tightlist
\item
  creazione/apertura e governo di progetti distinti
\item
  ownership project-local della documentazione e dei risultati
\end{itemize}

\subsubsection{What moved down}\label{what-moved-down-2}

Filesystem roots, templates, internal/external distinction, BAE types
and adapter boundaries become lower-level choices and requirements.

\begin{center}\rule{0.5\linewidth}{0.5pt}\end{center}

\subsection{MR-0005}\label{mr-0005}

\subsubsection{Current document}\label{current-document-4}

\paragraph{Original: MR-0005 - Common governed analysis
model}\label{original-mr-0005---common-governed-analysis-model}

\subparagraph{Intent}\label{intent-9}

Define the common model through which ThreatForge governs
methodology-specific analyses based on the Base Analysis and produces
repeatable derived representations.

\subparagraph{Context}\label{context-9}

The BAE registry is the methodology-neutral canonical source for the
analyzed system. Transformation of BAEs into methodology-specific
analysis elements involves expert judgment when the selected method does
not provide an unambiguous mapping. Reviewed and accepted analysis
registries provide the inputs for deterministic validation and derived
representation generation.

\subparagraph{Macro obligation}\label{macro-obligation-4}

\begin{itemize}
\tightlist
\item
  The common analysis model must preserve the BAE registry as the
  methodology-neutral canonical source.
\item
  Each analysis must record the expert-approved methodological
  interpretation separately.
\item
  Methodology-specific registries must use the controlled taxonomies
  defined by their analysis method.
\item
  The common analysis model must distinguish expert judgment from
  subsequent deterministic transformations.
\item
  The DFD of an analysis must be derived deterministically from the
  governed registries accepted for that analysis.
\item
  The renderer must consume a derived projection.
\item
  The renderer must not contain Base Analysis or methodology-specific
  rules.
\item
  Method-specific taxonomies and rules must belong to dedicated
  Macro-requirements.
\item
  An analysis that identifies missing Base Analysis knowledge must
  record the discrepancy explicitly.
\item
  The analysis process must not modify BAE registry records
  automatically.
\end{itemize}

\subparagraph{Scope}\label{scope-5}

\begin{itemize}
\tightlist
\item
  Includes: common model and boundaries of governed analyses
\item
  Includes: relationship between the canonical Base Analysis and
  methodology-specific registries
\item
  Includes: separation between expert interpretation and deterministic
  transformations
\item
  Includes: deterministic DFD derivation from accepted analysis
  registries
\item
  Includes: separation between DFD projection and rendering
\item
  Excludes: STRIDE-specific categories and rules
\item
  Excludes: STRIDE-AI-specific categories and rules
\item
  Excludes: detailed definition of DFD detail levels
\item
  Excludes: automatic promotion of analysis discoveries into the BAE
  registry
\item
  Excludes: automatic replacement of expert judgment
\end{itemize}

\subparagraph{Non-goals}\label{non-goals-4}

\begin{itemize}
\tightlist
\item
  Define the complete registry schemas for every methodology
\item
  Define detailed diagram layout rules
\item
  Implement the complete correction lifecycle for elements missing from
  the Base Analysis
\end{itemize}

\subsubsection{Proposed rewrite for metamodel
test}\label{proposed-rewrite-for-metamodel-test-4}

\paragraph{Proposed: MR-0005 - Analisi governate e
ripetibili}\label{proposed-mr-0005---analisi-governate-e-ripetibili}

\subparagraph{Intent}\label{intent-10}

Consentire di applicare analisi differenti allo stesso progetto
mantenendo espliciti il giudizio esperto, la provenienza dei risultati e
le parti che possono essere ripetute deterministicamente.

\subparagraph{Context}\label{context-10}

Uno stesso progetto puo essere osservato con paradigmi e metodi
differenti. ThreatForge deve permettere queste interpretazioni senza
alterare implicitamente i fatti comuni del progetto e senza confondere
risultati metodologici con conoscenza di base.

\subparagraph{Stakeholders}\label{stakeholders-5}

\begin{itemize}
\tightlist
\item
  analisti
\item
  reviewer
\item
  responsabili di progetto
\item
  specialisti delle discipline di analisi
\item
  utenti che confrontano o rieseguono analisi
\end{itemize}

\subparagraph{Objective}\label{objective-5}

\begin{itemize}
\tightlist
\item
  Governare analisi metodologicamente differenti sopra una base comune
  del progetto.
\item
  Rendere tracciabili le interpretazioni e i risultati prodotti da
  ciascuna analisi.
\item
  Permettere di ripetere in modo deterministico le trasformazioni che
  non richiedono nuovo giudizio esperto.
\end{itemize}

\subparagraph{Scope (candidate semantics under
test)}\label{scope-candidate-semantics-under-test-4}

\begin{itemize}
\tightlist
\item
  ciclo comune delle analisi e loro tracciabilita
\item
  separazione tra conoscenza comune, interpretazione metodologica e
  risultati derivati
\end{itemize}

\subsubsection{What moved down in the first
reformulation}\label{what-moved-down-in-the-first-reformulation-1}

DFD, renderer, method taxonomies, plugin specifics and exact registry
shapes are lower-level analysis/model decisions, not MR obligations.

\subsubsection{Second reformulation after ADR-family
review}\label{second-reformulation-after-adr-family-review-1}

\paragraph{Current candidate: MR-0005 - Governed analysis of the
system}\label{current-candidate-mr-0005---governed-analysis-of-the-system}

\subparagraph{Intent}\label{intent-11}

Enable different analysis methods to examine governed project knowledge
and produce reviewable, traceable and repeatable results without
changing the underlying description of the system or hiding
method-specific interpretation.

\subparagraph{Context}\label{context-11}

A governed description of the system provides shared knowledge about
what is being developed, but different analysis methods can interpret
that knowledge from different perspectives and may require expert
judgment.

Without a governed analysis boundary, analytical interpretations may be
confused with project facts, results from different methods may become
difficult to compare or trace, and conclusions may lose their connection
to the project knowledge and requirements that motivated them.

\subparagraph{Stakeholders}\label{stakeholders-6}

\begin{itemize}
\tightlist
\item
  Analysts
\item
  Security and assurance specialists
\item
  Project owners and maintainers
\item
  Architects and developers
\item
  Reviewers responsible for evaluating analysis results
\end{itemize}

\subparagraph{Objective}\label{objective-6}

Allow one governed description of a system to support multiple analysis
methods while preserving the method used, the analyst's interpretation,
the evidence considered and the results produced.

\subparagraph{Scope}\label{scope-6}

Includes governed applications of analysis methods to project knowledge.

Includes explicit analytical interpretation and expert judgment.

Includes traceability from analysis results to the project knowledge and
requirements they concern.

Includes review of analytical results and feedback when analysis reveals
missing, ambiguous or inconsistent project knowledge.

Excludes modification of governed project knowledge merely because an
analysis method interprets it differently.

Excludes the internal taxonomy, rules or procedure of any specific
analysis methodology from the common analysis model.

\subparagraph{Assumptions}\label{assumptions-1}

Different analysis methods may interpret the same project knowledge
differently.

Some analysis activities require expert judgment and cannot be reduced
entirely to deterministic transformation.

\subparagraph{Constraints}\label{constraints-1}

Method-specific interpretation must remain distinguishable from governed
project knowledge.

Analysis results must retain enough provenance to identify the method,
relevant project knowledge and supporting analytical evidence.

Discoveries about missing or incorrect project knowledge must return
through the governed documentation process rather than silently changing
canonical project knowledge.

\subsubsection{Why the first reformulation was
reopened}\label{why-the-first-reformulation-was-reopened-1}

The first reformulation described MR-0005 as applying different analyses
over a common project base. That was correct but too close to the first
MR-0003 rewrite, which had also been framed around enabling analysis.

Inspection of the historical MR-0005 Decisions showed a coherent second
concern: expert analysis records, common findings and documentation
feedback, methodology-specific extension boundaries, and governed use of
accepted findings.

The second reformulation therefore centers MR-0005 on a different
question: \textbf{what do we learn by analysing the governed system
knowledge, and how do we govern that interpretation and its results?}

\begin{center}\rule{0.5\linewidth}{0.5pt}\end{center}

\subsection{Lesson learned from MR-0003 versus
MR-0005}\label{lesson-learned-from-mr-0003-versus-mr-0005}

\subsubsection{Observation}\label{observation}

The first rewrites made MR-0003 and MR-0005 look almost duplicative. A
merge initially appeared plausible. That conclusion was deliberately
suspended and the historical Decision families were inspected before
changing the MR decomposition.

\subsubsection{Evidence that preserved the distinction for further
testing}\label{evidence-that-preserved-the-distinction-for-further-testing}

\begin{itemize}
\tightlist
\item
  MR-0003 Decisions form a coherent family around \textbf{governed
  system knowledge}: representation, identity/provenance, source
  continuity, documentary extraction/review and the semantic boundary
  between facts and method-specific interpretation.
\item
  MR-0005 Decisions form a coherent family around \textbf{governed
  analytical interpretation and results}: expert analysis records,
  findings/feedback, methodology extension boundaries and eligibility of
  accepted results for downstream use.
\item
  The distinction survives removal of current solution names such as
  BAE, Analysis Record, Finding, DFD and plugin.
\item
  A dependency can be described without hierarchy: governed analysis
  consumes governed system knowledge, but consuming it does not make
  analysis the same concern as representing it.
\end{itemize}

\subsubsection{Questions that must be asked before merging or splitting
Macro
Requirements}\label{questions-that-must-be-asked-before-merging-or-splitting-macro-requirements}

\begin{enumerate}
\def\labelenumi{\arabic{enumi}.}
\tightlist
\item
  \textbf{Independent value:} does each candidate MR express a
  project-level result that remains meaningful without describing the
  other?
\item
  \textbf{Decision-family coherence:} do the descendant Decisions form
  distinct, internally coherent families of choices?
\item
  \textbf{Solution-vocabulary removal:} does the distinction remain
  after removing names of current components, schemas, models, tools and
  other chosen mechanisms?
\item
  \textbf{Input versus result:} is one concern about establishing
  governed knowledge while another is about interpreting or transforming
  that knowledge? If so, is that distinction meaningful to the project
  rather than merely an implementation pipeline?
\item
  \textbf{Dependency versus containment:} can the relation be expressed
  as \texttt{dependsOn} rather than by merging the concerns or making
  one a child of the other?
\item
  \textbf{Stakeholder explanation:} can a stakeholder competent in the
  project domain explain the difference using title + Intent without
  knowing the implementation?
\item
  \textbf{Merge stress test:} if the MRs were merged, would the
  resulting MR need to explain two substantially different families of
  Decisions or two different kinds of project value?
\item
  \textbf{Split stress test:} if they remain separate, is the separation
  justified by project concerns, or only by the fact that the current
  architecture has two layers/subsystems?
\end{enumerate}

\subsubsection{Current disposition}\label{current-disposition}

\textbf{No merge decision is taken.} The current working hypothesis
keeps MR-0003 and MR-0005 distinct because the historical Decision
families reveal two different semantic concerns. This remains a
case-study result, not a universal rule that every project must separate
system representation from analysis.

\subsubsection{General lesson for the
metamodel}\label{general-lesson-for-the-metamodel}

When refactoring an existing governed corpus, similarity between
rewritten MR prose is evidence to investigate, not sufficient evidence
to merge. The candidate MR boundary must also be tested against the
semantic responsibility of its descendant Decisions and against a
version of the concern stripped of current solution vocabulary.

\begin{center}\rule{0.5\linewidth}{0.5pt}\end{center}

\end{document}
